
\section{2. Related Work}

\subsection{2.1 Window-Based Sparse Attention}

\textbf{Longformer} (Beltagy et al., 2020): Sliding window of size $w$ + global 
tokens. Achieves $\mathcal{O}(L w)$ complexity with fixed $w$.

\textbf{BigBird} (Zaheer et al., 2020): Sliding window + random + global. 
Theoretical universality but still fixed patterns.

\textbf{Mistral} (Jiang et al., 2023): Sliding window of 4096 tokens. Modern 
decoder-only adoption.

\textbf{Limitation.} All use fixed window size for all heads.

\subsection{2.2 Learned Sparse Patterns}

\textbf{Adaptive Span} (Sukhbaatar et al., 2019): \textbf{Our direct predecessor.} 
Learns per-head span via continuous soft ramp. The original paper applied 
it to language modeling (Wikitext-103) with a specific training procedure 
(span penalty with alternating attention).

\textbf{Learned Sparsity Patterns} (Tay et al., 2021): Differentiable stride 
and dilation. Complements ASA-Span.

\textbf{SparseBERT} (Xia et al., 2022): Learned sparsity via distillation.

\textbf{Our work} generalizes Sukhbaatar et al. to modern architectures (BERT, 
RoBERTa, GPT-2) with modern tools (RoPE, GQA, LoRA).

\subsection{2.3 Content-Based Sparsity}

\textbf{Routing Transformer} (Roy et al., 2021): K-means clustering of queries 
and keys. Content-based but not distance-based.

\textbf{Reformer} (Kitaev et al., 2020): LSH-based bucketing. Content-based.

\textbf{Top-K Attention} (Gupta et al., 2021): Hard top-$k$ selection.

\textbf{Our work} is \textbf{distance-based}, which is orthogonal to content-based 
methods and can be combined.

\subsection{2.4 Differentiable Masking}

\textbf{Gumbel-Softmax} (Jang et al., 2017): Stochastic relaxation of discrete 
sampling.

\textbf{Sigmoid soft mask} (Our Adaptive Sparse Attention): Differentiable 
top-$k$ approximation.

\textbf{Soft ramp} (Our ASA-Span): Piecewise-linear differentiable window.

\subsection{2.5 Position of Our Work}

\begin{table}[H]
\centering
\begin{tabular}{lll}
\toprule
\textbf{Axis} & \textbf{Prior} & \textbf{ASA-Span} \\
\midrule
Window type & Fixed & **Learned per-head** \\
Differentiability & -- & **Soft ramp** \\
Regularization & -- & **L1 span penalty** \\
Distance-based & Yes & **Yes** \\
Modern architectures & Limited & **BERT, RoBERTa, GPT-2, T5** \\
RoPE support & No & **Yes** \\
GQA support & No & **Yes** \\
\bottomrule
\end{tabular}
\end{table}

\textbf{Key distinction.} ASA-Span is the first \textbf{modern} implementation of 
learned per-head spans with support for RoPE, GQA, and T5-style biases. It 
generalizes Sukhbaatar et al. (2019) to the current transformer ecosystem.

\subsection{2.6 Comparison Table}
